\documentclass[12pt,A4,A4pt]{article}

\usepackage{amsmath}
\renewcommand*\familydefault{\sfdefault} 
\usepackage[scaled]{helvet}
\renewcommand\familydefault{\sfdefault} 
\usepackage[T1]{fontenc}
%\usepackage[T1]{fontenc}
%\usepackage{arial} 
%\pagestyle{myheadings}{}
 %\markboth{left head}{right head}
\usepackage{amssymb}
%\usepackage{epsfig}
\usepackage{fancyhdr}
\usepackage{float}
%\usepackage[authoryear,round]{natbib}
\usepackage{apacite}
%\usepackage[sort&compress]{natbib}
\usepackage[pdftex]{graphicx}
\usepackage{color}
\usepackage[brazil]{babel}
\usepackage[utf8]{inputenc}
\usepackage{amsfonts}
\usepackage{textcomp}
\usepackage{tabularx}
\usepackage{multirow}
\usepackage{times}
\usepackage{multicol}
\usepackage{lscape}
\usepackage{subfig}
\usepackage{epstopdf}
\usepackage{setspace}
\usepackage{fancyhdr}
\usepackage{epigraph}
\usepackage{scrextend}
\usepackage[shortlabels]{enumitem}
\usepackage{ragged2e}


%\let\OLDthebibliography\thebibliography
%\renewcommand\thebibliography[1]{\OLDthebibliography{#1} \setlength{\parskip}{0pt}\setlength{\itemsep}{0pt plus 0.3ex}}

%\renewcommand{\chaptermark}[1]{\markboth{\MakeUppercase{#1}}{}}
%\renewcommand{\sectionmark}[1]{\markright{\MakeUppercase{#1}}{}}
\pagestyle{fancy}
\fancyhf{}
\renewcommand{\headrulewidth}{0cm}
%\renewcommand{\footrulewidth}{0}
\usepackage{pdfpages}
\usepackage{stackrel}
\usepackage[left=3cm, right=3cm, top=2.5cm, bottom=2.5cm]{geometry}
\newtheorem{exemplo}{Exemplo}
\newtheorem{definicao}{Definição}
\newtheorem{teorema}{Teorema}
\cfoot{{\fontsize{10pt}{\baselineskip}\selectfont \textbf{8\textsuperscript{o} MCSul / VIII SEMENGO - Universidade Federal do Rio Grande}}}



%%%%%
%tira os numeros da frente do autor
\makeatletter 
\renewcommand\@biblabel[1]{}
\makeatother
%%


\begin{document}
\pretolerance10000


\begin{figure}
\centering
\vspace{-1cm}
\begin{minipage}[c]{\textwidth}
\centering
    \includegraphics[width=6.2in]{cabecalho.eps}
\end{minipage}
\end{figure}



\begin{center}
\fontsize{16pt}{\baselineskip}\selectfont 
\textbf{{ANÁLISE DO ALGORITMO DE EVOLUÇÃO DIFERENCIAL APLICADO AO DESIGN CONSTRUCTAL DE UMA CAVIDADE EM FORMA DE DUPLO-T}}
\end{center}
\vspace{-0.9cm}
% \begin{center}
% \vspace{0.5cm}
% \textbf{\large{\textit{TÍTULO DO TRABALHO EM INGLÊS}}}
% \end{center}

\begin{flushright}
Gill V. Gonzales\footnote{Mestre, Instituto Federal Sul-Rio-Grandense - Campus Santana do Livramento, RS, Brazil, gillgonzales@ifsul.edu.br.}$^{,2}$

Liércio A. Isold\footnote{Doutor, Universidade Federal do Rio Grande - PPG em Modelagem Computacional, Rio Grande, RS, Brazil, liercioisoldi@furg.br.}

Luiz A. Oliveira Rocha\footnote{Doutor, Universidade do Vale do Rio dos Sinos - PPG em Engenharia Mecânica – São Leopoldo, RS, Brasil luizor@unisinos.br.}

Elizaldo D. dos Santo\footnote{Doutor, Universidade Federal do Rio Grande - PPG em Modelagem Computacional, Rio Grande, RS, Brazil, elizaldosantos@furg.br.}

Antônio J. Silva Neto\footnote{Doutor, Universidade do Estado do Rio de Janeiro, Instituto Politécnico, Nova Friburgo, RJ, Brazil, ajsneto@iprj.uerj.br.}

\end{flushright}

\begin{flushleft}
{\small \setstretch{0.5} \justify
\textbf{Resumo:} Neste trabalho é analisado o desempenho do algoritmo de Evolução Diferencial aplicado à otimização geométrica de um problema de transferência de calor. O problema consiste em uma cavidade em forma de Duplo-T inserida em um sólido com geração de calor constante e uniforme. As laterias do sólido estão em condições adiabáticas e o calor gerado só pode ser removido pela cavidade a qual é mantida a uma temperatura prescrita. Portanto, a geometria da cavidade influência diretamente a performance térmica do problema. A definição dos graus de liberdade e restrições do problema é realizada através do método de Constructal Design. A otimização geométrica é realizada através do algoritmo de Evolução Diferencial. Neste estudo, são avaliados os parâmetros do algoritmo, tais como o operador de mutação, taxa de cruzamento, fator de amplificação do cálculo diferencial e número de iterações ou avaliação da função objetivo. O número de iteração está relacionado com o número de gerações e de indivíduos na população. O objetivo principal do trabalho é avaliar os parâmetros de configuração do algoritmo de Evolução Diferencial e qual a influência na reprodução dos efeitos dos graus de liberdade avaliados sobre a geometria ótima e performance térmica do problema investigado.
%192/300

\vspace{0.3cm}

\noindent\textbf{Palavras-chave:} Otimização Geométrica. Transferência de Calor. Palavra-chave. Palavra-chave. Palavra-chave.}
\end{flushleft}


% \noindent\textbf{Abstract:} O resumo em inglês apresenta-se logo após o resumo em português. Sugere-se encaminhar o texto para ser traduzido por um profissional. Caso o artigo esteja em inglês, deve haver uma versão do título, resumo e palavras-chave em português ou espanhol.
% \vspace{0.3cm}

% \noindent\textbf{Keywords:} Keywords. Keywords. Keywords. Keywords. Keywords.}

\newpage
 \onehalfspacing
\section{Introdução}

 {\fontsize{12pt}{\baselineskip}\selectfont}
 
\hspace{0.5cm}O artigo não precisa, necessariamente, contar com títulos que dividam ou subdividam sua estrutura, embora seja recomendado. Caso os autores optem por utilizar títulos e subtítulos, eles devem ser numerados desde a introdução.

A introdução é o texto de abertura do artigo. Deve mostrar a relevância do problema científico que motivou a realização do trabalho para área, apresentar a finalidade e os objetivos do trabalho proposto, destacando o aspecto inédito do trabalho e sua contribuição para área, bem como outras questões relevantes que explorem a temática abordada. Portanto, segue uma estrutura que permita delinear o assunto foco do trabalho, as referências que dão embasamento à discussão do tema, a justificativa para a escolha do tema e da importância do trabalho, o problema, os objetivos e o método utilizado. 

Quanto à formatação do corpo de texto emprega-se fonte 12 Arial, com espaçamento 1,5 entrelinhas e alinhamento justificado. 

Os anais do 8\textsuperscript{o} MCSUL e VIII SEMENGO serão publicados em formato digital (PDF - Portable Document Format). Os artigos devem ser formatados estritamente de acordo com as instruções aqui fornecidas. 

Artigos completos devem conter até $15$ páginas, incluindo Tabelas e Figuras. Os arquivos não devem ultrapassar 2MB.


\section{Seções}
\label{secoes}

\hspace{0.5cm}Seções e seus títulos são declaradas na forma usual do \LaTeX a partir dos seguintes comandos:

\begin{verbatim}
    \section
    \subsection
    \subsubsection
    \paragraph
\end{verbatim}

\section{Citações}
\hspace{0.5cm}As citações seguem o padrão APA\footnote{http://linorg.usp.br/CTAN/biblio/bibtex/contrib/apacite/apacite.pdf} (\textit{American Psychological Association}). Estas podem ser feitas pela utilização do comando \verb|\cite{}| como usual. O modelo irá produzir a citação na forma (Autor, Ano). Para produzir a citação na forma Autor(Ano) o comando  \verb|\citeA{}| deve usado.

Assim os números de cada citação são automaticamente organizados conforme o estilo do modelo. Por exemplo:

\begin{verbatim}
    "... deformation with a lateral resolution of the order of 100 nm
    \cite{Mizuuchi2005,Nakanishi2007}.  The recent development of 
    compact coherent EUV laser sources, \citeA{Sch2000} has opened ..."
\end{verbatim}
gera o texto da seguinte forma:

... deformation with a lateral resolution of the order of 100 nm \cite{Nakanishi2007},\break \cite{Mizuuchi2002}. The recent development of compact coherent EUV laser sources, \citeA{scht2000} has opened ...

Texto exemplo para citações.

Using synchrotron radiation Gabor and Fourier holograms have been demonstrated \cite{Mizuuchi2003}
 with spatial resolution below 100 nm at SXR wavelengths.  Compact EUV sources based on high harmonic generation (HHG) were also used to demonstrate table-top in-line EUV holography with a spatial resolution of 7.9 m  and 0.8 m.  Time resolved holographic imaging, that exploits the short pulsewidth of the HHG sources, was also implemented to study the ultrafast dynamics of surface deformation with a lateral resolution of the order of 100 nm \cite{Mizuuchi2005,Nakanishi2007}.  The recent development of compact coherent EUV laser sources, \citeA{scht2000} has opened new opportunities for the implementation of novel imaging schemes with nanometer-scale resolution that fit on a table-top \cite{Mizuuchi2003,scht2000}. In this paper, we present a proof of principle experiment in which we demonstrate that three dimensional imaging in a volume may be obtained from a single high numerical  aperture (NA) hologram obtained with a table top EUV laser. Gabor holograms were numerically reconstructed over a range of image planes by sweeping the propagation distance. This numerical \nobreak{sectioning} technique for holography is verified to produce a robust three dimension image of a test object.  

\section{Equações}
As equações são criadas pela utilização do ambiente tradicional \LaTeX:

\begin{verbatim}
    \begin{equation}
        \label{eqn_exemplo}
        x = \sum\limits_{i=0}^{z} 2^{i}Q
    \end{equation}
\end{verbatim}
o que produz:

\begin{equation}
     \label{eqn_exemplo}
      x = \sum\limits_{i=0}^{z} 2^{i}Q
\end{equation}

Utilize o ambiente \verb|displaymath| caso não deseje numeração de alguma equação. Quando alguma referência for feita a qualquer equação esta deve fazer a indexação no próprio texto pelo comando \verb|\ref{}| entre parênteses, precedido da sigla "Eq.", conforme a seguir:

\begin{verbatim}
    ... de acordo com a Eq.(\ref{eqn_exemplo}).
\end{verbatim}
o que produz:

"... de acordo com a Eq.(\ref{eqn_exemplo})."


\subsection{Equações Multilinhas}

O uso de equações multilinhas é exemplificado a seguir:

\begin{verbatim}
\setlength{\arraycolsep}{0.0em}
  \begin{eqnarray}
  Z&{}={}&x_1 + x_2 + x_3 + x_4 + x_5 + x_6  \nonumber \\
  &&+a + b\\
  &&+{}a + b\\
  &&{}+a + b\\
  &&{+}\:a + b
  \end{eqnarray}
\setlength{\arraycolsep}{5pt}
\end{verbatim}

O que produz:

\setlength{\arraycolsep}{0.0em}
\begin{eqnarray}
Z&{}={}&x_1 + x_2 + x_3 + x_4 + x_5 + x_6  \nonumber \\
&&+a + b\\
&&+{}a + b\\
&&{}+a + b\\
&&{+}\:a + b
\end{eqnarray}
\setlength{\arraycolsep}{5pt}

\subsection{Outras formas de Equações}

As equações ou variáveis matemáticas podem ser utilizadas em linha de texto. Para tanto a sintaxe se apresenta da seguinte forma:

\begin{verbatim}
"Sendo $X_1, X_2, \ldots, X_n$ uma sequência de variáveis independentes e
identicamente distribuídas aleatóriamente com $\text{E}[X_i] = \mu$ e
$\text{Var}[X_i] = \sigma^2 < \infty$, e sendo..."
\end{verbatim}

A qual produz o parágrafo a seguir:\\


"Sendo $X_1, X_2, \ldots, X_n$ uma sequência de variáveis independentes e identicamente distribuídas aleatóriamente com $\text{E}[X_i] = \mu$ e $\text{Var}[X_i] = \sigma^2 < \infty$, e sendo
\begin{equation}
S_n = \frac{X_1 + X_2 + \cdots + X_n}{n}
      = \frac{1}{n}\sum_{i}^{n} X_i
\label{eq:refname1}
\end{equation}

denotam seu significado. Então, quando $n$ tende ao infinito, as veriáveis aleatórias $\sqrt{n}(S_n - \mu)$ convergem em uma distribuição normal $\mathcal{N}(0, \sigma^2)$. Isto é apresentado na Eq.(\ref{eq:refname1})"\\



\begin{equation}
 \frac{\partial T}{\partial t} =\kappa \big(\frac{\partial^2T}{\partial x^2}+\frac{\partial^2T}{\partial y^2} \big) 
 \label{eq:refname2}
\end{equation}

Ou mesmo a Eq.(\ref{eq:refname3}):

\begin{equation}
  \vec{F}=\mathbf{m} \frac{d \vec{v}}{dt} + \vec{v}\frac{dm}{dt} 
  \label{eq:refname3}
\end{equation}

\begin{subequations}
\begin{equation}
  \operatorname{min}_{a,b,c} 
  \frac{1}{2}\mathbf{w}^{T}\mathbf{w} + C \sum_{i=1}^{l}\xi_{i}
  \label{eq:refname4a}
\end{equation}    
\begin{equation}
  y_{i}\left(\mathbf{w}^{T}\phi(x_{i})+b\right)
  \label{eq:refname4b}
\end{equation}
\end{subequations}

Podendo usar subequações para um conjunto de equações do mesmo grupo, como a Eq.(\ref{eq:refname4a}) e (\ref{eq:refname4b}), por exemplo.

A sintaxe para produzir o conjunto de Equações anterior é dada por:

\begin{verbatim}

\begin{equation}
\frac{\partial T}{\partial t} =\kappa \big(\frac{\partial^2T}
{\partial x^2}+\frac{\partial^2T}{\partial y^2} \big) 
 \label{eq:refname2}
\end{equation}

Ou mesmo a Eq.(\ref{eq:refname3}):

\begin{equation}
  \vec{F}=\mathbf{m} \frac{d \vec{v}}{dt} + \vec{v}\frac{dm}{dt} 
  \label{eq:refname3}
\end{equation}

\begin{subequations}
\begin{equation}
  \operatorname{min}_{a,b,c} 
  \frac{1}{2}\mathbf{w}^{T}\mathbf{w} + C \sum_{i=1}^{l}\xi_{i}
  \label{eq:refname4a}
\end{equation}    
\begin{equation}
  y_{i}\left(\mathbf{w}^{T}\phi(x_{i})+b\right)
  \label{eq:refname4b}
\end{equation}
\end{subequations}

\end{verbatim}

Lembrando que se você \emph{nomear} todas suas equações, a indexação fica muito mais fácil no momento em que as mesmas precisam ser citadas.
Outra observação importante: \textit{\textbf{as linhas de comandos desde quando se inicia a inserção de uma equação até quando se finaliza, NÃO PODEM APRESENTAR LINHAS EM BRANCO DE UMA LINHA DE COMANDO PARA A OUTRA.}}

\section{Estruturas Flutuantes}
\subsection{Figuras}
Os autores devem ter em mente que todo tipo de estrutura flutuante, seja figura ou tabela deve estar o mais próximo de sua referência no texto possível. Em \LaTeX pode-se dar prioridades diferentes no posicionamento destas, através das letras \textit{h, t, b, p}, as quais signficam \textit{here} (aqui), \textit{top} (topo), \textit{bottom} (base) e \textit{page} (página). Estas letras deve ser colocadas imediatamente ao lado do comando de inserção do elemento flutuante. O padrão de formatação \LaTeX automaticamente tenta colocar estes elementos de forma a deixá-los próximos de suas referências no texto. Estes comando são utilizados para casos em que o autor não esteja satisfeito com o posicionamento automático.

A inserção de uma figura, por exemplo, no "topo" da página como posição preferencial é feita a partir dos seguintes comandos (usuais em \LaTeX). A Fig.(\ref{fig_env1}) é um exemplo de inserção.

\begin{verbatim}
    
    \begin{figure}[t]
        \centering
        \includegraphics[width=21pc]{mouse}
        \caption{(a) Diagrama do aparato experimental.}
        \label{fig_env1}
    \end{figure}
    
\end{verbatim}

O que produz a figura no artigo. Detalhes como: ajustar por largura, ou por escala, ou por medidas específicas estão na documentação e manual do \LaTeX.

\begin{figure}[t]
\centering
\includegraphics[width=21pc]{mouse}
\caption{\fontsize{10pt}{\baselineskip}\selectfont Diagrama do aparato experimental.}
\label{fig_env1}
\end{figure}

As figuras podem ser inseridas lado a lado, dependendo da organização das informações.

Os seguintes comandos exemplificam a organização das figuras \textit{Lado a Lado} com única legenda:

\begin{verbatim}
    
\begin{figure}[htbp]
\centering
\includegraphics[scale=0.6]{mouse}
\quad
\includegraphics[scale=0.2]{mouse}
\caption{Firuras lado-a-lado com mesma legenda}
\label{fig:lado_a_lado}
\end{figure}

\end{verbatim}

Os comandos anteriores produzem a Fig.(\ref{fig:lado_a_lado}).

\begin{figure}[htbp]
\centering
\includegraphics[scale=0.8]{mouse}
\quad
\includegraphics[scale=0.5]{mouse}
\caption{\fontsize{10pt}{\baselineskip}\selectfont Figuras lado-a-lado com mesma legenda}
\label{fig:lado_a_lado}
\end{figure}


Ainda, figuras com legendas diferentes também são uma opção, podendo ser colocadas a partir dos seguintes comandos, conforme Fig.(\ref{fig:lado_a_lado}).

A referência cruzada das figuras no texto é feita a partir da referência das mesmas através do comando \verb|\ref{nome do rótulo}|. O rótulo da Figura é definido pelo comando \verb|\label{}|, conforme pode ser observado na sequência de comandos anteriores, para a inserção da figura.

\subsection{Tabelas}

As tabelas podem ser trabalhadas de maneira similar às figuras, apenas com algumas pequenas diferenças. Construir tabelas no ambiente \LaTeX é, a primeira vista, uma tarefa um tanto complexa. No entanto, existem sites e/ou suplementos para serem utilizados no \textit{Microsoft Excell}, por exemplo, os quais, a partir de uma tabela construída em um ambiente "amigável" pode ser "traduzida" para um código a ser usado em \LaTeX. Um editor de tabelas \emph{online} para \LaTeX\ é, por exemplo, o Tables Generator\footnote{http://www.tablesgenerator.com}.

Exemplos de geração de tabelas incluem o seguinte código:

\begin{verbatim}

\begin{table}[htbp]
\centering
\caption{Funções de Forma para Elementos Quadráticos}
\begin{tabular}{ccc}
\hline
local node & $\{N\}_m$ & $\{\Phi_i\}_m$ $(i=x,y,z)$ \\
\hline
$m = 1$ & $L_1(2L_1-1)$ & $\Phi_{i1}$ \\
$m = 2$ & $L_2(2L_2-1)$ & $\Phi_{i2}$ \\
$m = 3$ & $L_3=4L_1L_2$ & $\Phi_{i3}$ \\
\hline
\end{tabular}
  \label{tab:Funções-Forma}
\end{table}

\end{verbatim}

O resultado da sequência anteior é apresentado na Tab.\ref{tab:Funções-Forma}.

\begin{table}[htbp]
\small
\centering
\caption{\small Funções de Forma para Elementos Quadráticos}
\begin{tabular}{ccc}
\hline
local node & $\{N\}_m$ & $\{\Phi_i\}_m$ $(i=x,y,z)$ \\
\hline
$m = 1$ & $L_1(2L_1-1)$ & $\Phi_{i1}$ \\
$m = 2$ & $L_2(2L_2-1)$ & $\Phi_{i2}$ \\
$m = 3$ & $L_3=4L_1L_2$ & $\Phi_{i3}$ \\
\hline
\end{tabular}
  \label{tab:Funções-Forma}
\end{table}

A referência cruzada das tabelas no texto é feita a partir da citação das mesmas através do comando \verb|\ref{nome do rótulo}|. O rótulo da Tabela é definido pelo comando \verb|\label{}|, conforme pode ser observado na sequência de comandos anteriores, para a inserção da tabela.

\section{Listas}

\subsection{Itens}

As listas de itens podem ser colocadas através dos comandos \break \verb|\begin{itemize}| $ \ldots $ \verb|\end{itemize}|, ou \verb|\begin{enumerate}| $ \ldots $ \verb|\end{enumerate}|  como segue:



Para o caso em que se deseja utilizar a listagem por categoria alfabética automática, o comando conforme a seguir é utilizado (é necessário, para isto, incluir o pacote \break \verb|\usepackage[shortlabels]{enumitem}| no preâmbulo):

\begin{verbatim}
    
\begin{enumerate}[a)] 
\item Lorem ipsum dolor sit amet...
\item Praesent vitae nulla varius...
\item Praesent quis erat eleifend...
\item Mauris facilisis odio eu ...
\item Curabitur convallis odio...
\end{enumerate}

\end{verbatim}

tendo como resultado:

\begin{enumerate}[a)] 
\item Lorem ipsum dolor sit amet...
\item Praesent vitae nulla varius...
\item Praesent quis erat eleifend...
\item Mauris facilisis odio eu ...
\item Curabitur convallis odio...
\end{enumerate}

Ou em alíneas maiúsculas:

\begin{verbatim}
    
\begin{enumerate}[A)]
\item Lorem ipsum dolor sit amet...
\item Praesent vitae nulla varius...
\item Praesent quis erat eleifend...
\item Mauris facilisis odio eu ...
\item Curabitur convallis odio...
\end{enumerate}

\end{verbatim}

\begin{enumerate}[A)]   
\item Lorem ipsum dolor sit amet...
\item Praesent vitae nulla varius...
\item Praesent quis erat eleifend...
\item Mauris facilisis odio eu ...
\item Curabitur convallis odio...
\end{enumerate}
%\end{enumerate}

Ou ainda, em números romanos:

\begin{verbatim}
    
\begin{enumerate}[i.]
\item Lorem ipsum dolor sit amet...
\item Praesent vitae nulla varius...
\item Praesent quis erat eleifend...
\item Mauris facilisis odio eu ...
\item Curabitur convallis odio...
\end{enumerate}

\end{verbatim}

\begin{enumerate}[i.]
\item Lorem ipsum dolor sit amet...
\item Praesent vitae nulla varius...
\item Praesent quis erat eleifend...
\item Mauris facilisis odio eu ...
\item Curabitur convallis odio...
\end{enumerate}

E ainda, a opção sem nenhum indicador de sequência:

\begin{verbatim}
    
\begin{itemize} 
\item Lorem ipsum dolor sit amet...
\item Praesent vitae nulla varius...
\item Praesent quis erat eleifend...
\item Mauris facilisis odio eu ...
\item Curabitur convallis odio...
\end{itemize}

\end{verbatim}

\begin{itemize} 
\item Lorem ipsum dolor sit amet...
\item Praesent vitae nulla varius...
\item Praesent quis erat eleifend...
\item Mauris facilisis odio eu ...
\item Curabitur convallis odio...
\end{itemize}

%**********************************************************************************
\section*{Agradecimentos}
\hspace{0.5cm}Agradecimentos e outras seções não numeradas são criadas pela utilização de um asterisco posterior ao comando, \verb|\section*{Agradecimentos}|.

OBS: Parte dos textos utilizados neste modelo para exemplificar seções, quadros, tabelas, etc., foram gentilmente cedidos pela Revista Mundi\footnote{http://periodicos.ifpr.edu.br}, do IFPR.

%**********************************************************************************
\section*{Referencial bibliográfico}
\hspace{0.5cm}As referências bibliográficas são geradas automaticamente pela utilização do pacote \emph{apacite}. É uma seção não numerada conforme apresentado a seguir:


\bibliographystyle{apacite}
\bibliography{referencial}


\end{document}
